% from ega2-1-fr.tex
\oldpage[II]{76}

$Y$-morphisme $j\circ r$ correspond à
\[
  q^*(\mathcal{E})\xrightarrow{q^*(v)}q^*(\mathcal{F})
  \xrightarrow{\varphi}\mathcal{L}~;
\]
cela résulte encore de la définition donnée en
(\hyperref[II.4.2.1-fr]{4.2.1}).
\end{env}

\begin{env}[4.2.10]
\label{II.4.2.10-fr}
Soit $\psi:Y'\to Y$ un morphisme, et posons
$\mathcal{E}'=\psi^*(\mathcal{E})$. Si le $Y$-morphisme $r:X\to P$ correspond
à l'homomorphisme $\varphi:q^*(\mathcal{E})\to\mathcal{L}$, le
$Y'$-morphisme
\[
  r_{(Y')}:X_{(Y')}\longrightarrow P'=\mathbf{P}(\mathcal{E}')
\]
correspond à
$\varphi_{(Y')}:q_{(Y')}^*(\mathcal{E}')=
  q^*(\mathcal{E})\otimes_Y\mathcal{O}_{Y'}\to
\mathcal{L}\otimes_Y\mathcal{O}_{Y'}$. En effet, en vertu de
(\hyperref[II.4.1.3.1-fr]{4.1.3.1}), on a le diagramme commutatif
\[
  \xymatrix{
    Y' \ar[d]
    & P'=P_{(Y')} \ar[l]_{p_{(Y')}} \ar[d]^{u}
    & X_{(Y')} \ar[l]_{r_{(Y')}} \ar[d]^{v}
    \\Y
    & P \ar[l]_{p}
    & X \ar[l]_{r}
  }
\]
D'après (\hyperref[II.4.1.3.2-fr]{4.1.3.2}), on a
\[
  (r_{(Y')})^*(\mathcal{O}_{P'}(1))
  =(r_{(Y')})^*(u^*(\mathcal{O}_P(1)))
  =v^*(r^*(\mathcal{O}_P(1)))
  =v^*(\mathcal{L})=\mathcal{L}\otimes_Y\mathcal{O}_{Y'}~;
\]
d'autre part, $u^*(\alpha_1^\sharp)$ n'est autre que l'homomorphisme
canonique
$\alpha_1^\sharp:(p_{(Y')})^*(\mathcal{E}')\to\mathcal{O}_{P'}(1)$~; cela
se voit en explicitant les homomorphismes canoniques $\alpha_1^\sharp$ dans
$P$ et $P'$ comme dans (\hyperref[II.4.1.6-fr]{4.1.6}). D'où notre assertion.
\end{env}

\subsection{Le morphisme de Segre.}
\label{subsection:II.4.3-fr}

\begin{env}[4.3.1]
\label{II.4.3.1-fr}
Soient $Y$ un préschéma, $\mathcal{E}$, $\mathcal{F}$ deux
$\mathcal{O}_Y$-Modules quasi-cohérents~; posons
$P_1=\mathbf{P}(\mathcal{E})$, $P_2=\mathbf{P}(\mathcal{F})$, et désignons
par $p_1:P_1\to Y$, $p_2:P_2\to Y$ les morphismes structuraux. Soit
$Q=P_1\times_YP_2$, et soient $q_1:Q\to P_1$, $q_2:Q\to P_2$ les
projections canoniques~; le $\mathcal{O}_Q$-Module
\[
  \mathcal{L}=\mathcal{O}_{P_1}(1)\otimes_Y\mathcal{O}_{P_2}(1)
  :=q_1^*(\mathcal{O}_{P_1}(1))\otimes_{\mathcal{O}_Q}
  q_2^*(\mathcal{O}_{P_2}(1))
\]
est inversible comme produit tensoriel de deux $\mathcal{O}_Q$-Modules
inversibles (\hyperref[0.5.4.4-fr]{0, 5.4.4}). D'autre part, si
$r=p_1\circ q_1=p_2\circ q_2$ est le morphisme structural $Q\to Y$, on a
\[
  r^*(\mathcal{E}\otimes_{\mathcal{O}_Y}\mathcal{F})
  =q_1^*(p_1^*(\mathcal{E}))\otimes_{\mathcal{O}_Q}
  q_2^*(p_2^*(\mathcal{F}))
\]
(\hyperref[0.4.3.3-fr]{0, 4.3.3})~; les homomorphismes canoniques surjectifs
(\hyperref[II.4.1.5.1-fr]{4.1.5.1})
$p_1^*(\mathcal{E})\to\mathcal{O}_{P_1}(1)$,
$p_2^*(\mathcal{F})\to\mathcal{O}_{P_2}(1)$ donnent donc, par produit
tensoriel, un homomorphisme canonique
\[
\label{II.4.3.1.1-fr}
  s:r^*(\mathcal{E}\otimes_{\mathcal{O}_Y}\mathcal{F})\longrightarrow
  \mathcal{L}
\tag{4.3.1.1}
\]
évidemment surjectif~; on en déduit donc
(\hyperref[II.4.2.2-fr]{4.2.2}) un morphisme canonique, dit
\emph{morphisme de Segre}~:
\[
\label{II.4.3.1.2-fr}
  \varsigma:\mathbf{P}(\mathcal{E})\times_Y\mathbf{P}(\mathcal{F})
  \longrightarrow
  \mathbf{P}(\mathcal{E}\otimes_{\mathcal{O}_Y}\mathcal{F}).
\tag{4.3.1.2}
\]
Explicitons le morphisme $\varsigma$, lorsque $Y=\operatorname{Spec}(A)$ est
affine, $\mathcal{E}=\widetilde{E}$, $\mathcal{F}=\widetilde{F}$, $E$ et $F$
étant deux $A$-modules, d'où
$\mathcal{E}\otimes_{\mathcal{O}_Y}\mathcal{F}=\widetilde{E\otimes_A F}$
(\hyperref[I.1.3.12-fr]{I, 1.3.12})~; posons
$R=\mathbf{S}_A(E)$, $S=\mathbf{S}_A(F)$,
$T=\mathbf{S}_A(E\otimes_A F)$~; soient $f\in E$, $g\in F$, et considérons
l'ouvert affine
\[
  D_+(f)\times_YD_+(g)=\operatorname{Spec}(B)
\]

